% 2. Objetivos
\begin{frame}{Objetivo general}
  \begin{itemize}
    \item Construir una herramienta pedagógica sólida para entender y practicar el análisis sintáctico descendente predictivo (LL(1)), uniendo teoría y visualización paso a paso.
  \end{itemize}
\end{frame}

\begin{frame}{Objetivos específicos}
  \begin{itemize}
    \item Robustecer el flujo completo: transformaciones de gramáticas cuando proceda, cálculo correcto de Primero/Siguiente y construcción de la tabla predictiva LL(1) sin ambigüedades.
    \item Mejorar la experiencia de aprendizaje con un asistente claro y con derivación/árbol en tiempo real.
    \item Facilitar la práctica: edición de gramáticas, funciones de recuperación de errores y generación de informes.
  \end{itemize}
\end{frame}

\begin{frame}{Calidad del software}
  El proyecto persigue una mejora tangible en la calidad del software:
  \begin{itemize}
    \item \textbf{Consistencia}: datos y estado coherentes al navegar y al guardar.
    \item \textbf{Eficiencia}: cálculos y renderizados ágiles en equipos estándar.
    \item \textbf{Fiabilidad}: resultados deterministas y manejo sólido de errores.
    \item \textbf{Usabilidad}: interfaz limpia, asistente por pasos y mensajes claros.
    \item \textbf{Compatibilidad}: correcto funcionamiento con recursos y formatos esperados.
    \item \textbf{Portabilidad}: funcionamiento equivalente en Windows, macOS y Linux.
    \item \textbf{Mantenibilidad}: código modular, patrones conocidos y documentación.
    \item \textbf{Extensibilidad}: base preparada para nuevas funciones y vistas.
  \end{itemize}
\end{frame}

\begin{frame}{Resultados esperados}
  \begin{itemize}
    \item Experiencias de simulación consistentes, sin pérdidas de datos al navegar entre vistas.
    \item Portabilidad real (Windows, macOS y Linux) y una interfaz limpia basada en pestañas.
    \item Un soporte didáctico reutilizable en clase y para el autoestudio (PDF/HTML).
  \end{itemize}
\end{frame}


